\subsubsection{Corrections needed before importing the other claims}

\begin{enumerate}
\item \textbf{Uniform Fourier levels have additive influences.}
If $p=\sum_{k=0}^d p_k$ is its homogeneous Walsh decomposition, then
with our normalized flip convention,
\[
 I_i^U(p)=\sum_{S\ni i}\widehat p(S)^2
         =\sum_{k=1}^d I_i^U(p_k).
\]
The contrary sentence in the handoff is false under uniform input.
What a levelwise fcb argument must still track is the norm of the
selected level, including its ambient degree parameter. Under
correlated laws, Walsh orthogonality need not survive.

\item \textbf{The stated polynomial fcb target is already a uniform-AA
corollary.} The primary definition satisfies
$\|p\|_\infty\le\|p\|_{\mathrm{fcb},d}$; see
\href{https://arxiv.org/abs/2304.06713v2}{Escudero Guti\'errez,
Proposition 2.3}. Hence $\|p\|_{\mathrm{fcb},d}\le1$ is a stronger
restriction than supremum boundedness and puts $p$ within the
uniform-AA theorem imported here. The sharper degree-independent
$\operatorname{Var}^2$ estimate, or a proof exploiting fcb directly,
is a separate target. This supplies no correlated-Gibbs transfer.

\item \textbf{Repair the mean-square comparison witness.}
The handoff's $f=1+n^{-1}\sum_iY_i$, with independent mean-$1/2$
signs $Y_i$, gives a squared-norm ratio tending to $9/4$; that example
alone does not refute an arbitrary finite comparison constant. The
intended claim is proved by the bounded affine polynomial
$M=n^{-1}\sum_iY_i$:
\[
 \mathbb E_U M^2=1/n,\qquad
 \mathbb E_\mu M^2=\tfrac14+\tfrac{3}{4n},\qquad
 \frac{\mathbb E_\mu M^2}{\mathbb E_U M^2}=\frac{n+3}{4}.
\]
Independent fields themselves have full AA. This is a failure of a
comparison method, not of simulation.

\item \textbf{Do not use the reported inverse-atom prefactor.}
The handoff's project-level two-bit influence formula has a factor
$\delta_{\min}^{-6}$. It cannot hold with a positive prefactor uniform
in $\delta_{\min}$: take a nearly-copy two-bit law with fair marginals
and $p(x)=x_1$ on $\{0,1\}^2$. Its variance is $1/4$ and unnormalized
maximum flip influence is one, while the proposed right side diverges
as the least atom probability tends to zero. The source formula needs
its exact proof and parameter dependence; we do not silently repair
its exponent or rely on it. Our full-AA block theorem has its own proof.

\item \textbf{Boolean degree does not imply $d^2\ge n$.}
For the addressing function with $k$ address bits and $2^k$ data bits,
all $n=k+2^k$ variables are relevant, while the exact degree is $d=k+1$.
Expand it as $\sum_a\mathbf1[\text{address}=a]x_a$ to see the upper
bound; the top coefficient involving a fixed data bit is nonzero.
The handoff's blanket Boolean shortcut in Task 2 is therefore unavailable.

\item \textbf{Critical covariance needs a separate argument.}
On a free $b$-ary Ising tree of depth $h$,
$\operatorname{Cov}(X_o,X_v)=\tanh(\beta)^{\operatorname{dist}(o,v)}$.
At $\tanh\beta=1/b$, the root covariance row sum is $h+1$, not a
dimension-uniform constant. Thus that proof route cannot establish the
reported critical endpoint without another argument. This does not
refute a separate influence theorem, and our forest simulation does
not use covariance-row summability. On the square lattice the usual
field-uniform coefficient bound is $4\tanh\beta<1$, reaching
$\operatorname{arctanh}(1/4)$; $\beta<1/4$ is the more conservative
linear row-norm criterion. Neither should be conflated with a
convention-free exact physical or full-AA threshold.
\end{enumerate}
